% Bagian: first-order-logic
% Bab: axiomatic-deduction
% Subbagian: deduction-theorem-quantifiers

\documentclass[../../../include/open-logic-section]{subfiles}

\begin{document}

\olfileid[id]{fol}{axd}{ddq}

\olsection{Teorema Deduksi dengan Kuantor}

\begin{thm}[Teorema Deduksi]
\ollabel{thm:deduction-thm-q} Jika $\Gamma \cup \{!A\} \Proves !B$,
maka $\Gamma \Proves !A \lif !B$.
\end{thm}

\begin{proof}
Sekali lagi, kita melakukan induksi pada panjang !!{derivation} dari
$!B$ berdasarkan $\Gamma \cup \{!A\}$.

Pembuktian basis induksi identik dengan pembuktian basis induksi dalam
pembuktian~\olref[ded]{thm:deduction-thm}.

Untuk langkah induksi, andaikan sekali lagi bahwa !!{derivation} dari
$!B$ berdasarkan $\Gamma \cup \{!A\}$ berakhir dengan suatu langkah~$!B$
yang dijustifikasi oleh aturan inferensi. Jika aturan inferensinya adalah
modus ponens, kita melanjutkan seperti dalam pembuktian
\olref[ded]{thm:deduction-thm}. Jika aturan inferensinya adalah \QR, kita
mengetahui bahwa $!B \ident !C \lif \lforall[x][!D(x)]$ dan bahwa
!!a{formula} berbentuk $!C \lif !D(a)$ muncul lebih awal dalam
!!{derivation} tersebut, dengan $a$ tidak muncul dalam~$!C$, $!A$,
$!D(x)$, ataupun $\Gamma$. Dengan demikian, kita mempunyai
\begin{align*}
  \Gamma \cup \{!A\} & \Proves !C \lif !D(a),\\
  \intertext{dan hipotesis induksi berlaku; yakni, kita mempunyai}
    \Gamma & \Proves !A \lif (!C \lif !D(a)).\\
  \intertext{Dengan menggunakan}
  & \Proves (!A \lif (!C \lif !D(a))) \lif ((!A \land !C) \lif !D(a))\\
  \intertext{dan modus ponens, kita memperoleh}
  \Gamma & \Proves (!A \land !C) \lif !D(a).\\
  \intertext{Karena syarat variabel eigen tetap terpenuhi, kita dapat
    menambahkan satu langkah yang dijustifikasi oleh \QR{} pada
    !!{derivation} ini dan memperoleh}
    \Gamma & \Proves (!A \land !C) \lif \lforall[x][!D(x)].\\
    \intertext{Kita juga mempunyai}
    & \Proves ((!A \land !C) \lif \lforall[x][!D(x)]) \lif
      (!A \lif (!C \lif \lforall[x][!D(x)])),\\
    \intertext{sehingga, menurut modus ponens,}
    \Gamma & \Proves !A \lif (!C \lif \lforall[x][!D(x)]),
\end{align*}
yakni, $\Gamma \Proves !A \lif !B$.

Kasus ketika $!B$ dijustifikasi oleh aturan \QR, tetapi berbentuk
$\lexists[x][!D(x)] \lif !C$, kita tinggalkan sebagai latihan.
\end{proof}

\begin{prob}
  Lengkapi pembuktian \olref[fol][axd][ddq]{thm:deduction-thm-q}.
\end{prob}

\end{document}
